% Documento de requerimientos del proyecto.
% Compilar desde la raíz del repo con:  mise run docs:pdf
% (o, dentro de Documentacion/:  tectonic -X compile main.tex)
%
% Cada capítulo vive en su propia subcarpeta (NN-nombre/nombre.tex), para que
% cada parte del equipo pueda editar la suya sin pisar las demás. Si un capítulo
% necesita imágenes, van en su misma subcarpeta.

\documentclass[11pt,a4paper,oneside]{report}

\usepackage[spanish,es-noquoting,es-nodecimaldot]{babel}
\usepackage{fontspec}
\usepackage[margin=2.5cm]{geometry}
\usepackage{amsmath}
\usepackage{booktabs}
\usepackage{tabularx}
\usepackage{xltabular}
\usepackage{enumitem}
\usepackage{xcolor}
\usepackage{titlesec}
\usepackage[hidelinks]{hyperref}
\usepackage{xurl}

\setlist{nosep, leftmargin=*}
\setlength{\parskip}{0.5em}
\setlength{\parindent}{0pt}

% Capítulos con título compacto (sin "Capítulo N" en una página aparte)
\titleformat{\chapter}[hang]{\Huge\bfseries}{\thechapter}{0.6em}{}
\titlespacing*{\chapter}{0pt}{0pt}{1.5em}

% ---- Convenciones -----------------------------------------------------------
% Prioridad de cada requerimiento (MoSCoW simplificado)
\newcommand{\imp}{\textcolor{red!70!black}{\textbf{Imprescindible}}}
\newcommand{\rel}{\textcolor{orange!80!black}{\textbf{Importante}}}
\newcommand{\des}{\textcolor{gray}{\textbf{Deseable}}}
% Notas
\newcommand{\decision}[1]{\textcolor{teal}{\textbf{Decisión:}} #1}
\newcommand{\abierta}[1]{\textcolor{orange!80!black}{\textbf{Abierta:}} #1}

% Tabla de requerimientos: \begin{requerimientos} ID & texto & prioridad \\ ... \end{requerimientos}
\newenvironment{requerimientos}
  {\xltabular{\linewidth}{@{}l X l@{}}
    \toprule ID & Requerimiento & Prioridad \\ \midrule \endhead
    \bottomrule \endfoot}
  {\endxltabular}

\title{Documento de requerimientos\\[0.4em]
  \Large Calendario armonizado de focos de calor MODIS--VIIRS}
\author{Equipo Hots-Spots --- NASA Space Apps Challenge 2026}
\date{14 de noviembre de 2026 \quad·\quad Versión 0.1 (borrador)}

\begin{document}
\maketitle
\tableofcontents

\chapter{Introducción}\label{cap:introduccion}

\section{Propósito del documento}

Fijar \emph{qué} tiene que hacer el proyecto antes de seguir construyéndolo: los
requerimientos de los datos, de la aplicación web y de calidad, y la estructura con
la que se organiza el trabajo. Es la referencia común del equipo: si algo no está
aquí, se discute y se añade; si cambia, se cambia aquí primero.

Documentos relacionados:

\begin{tabularx}{\linewidth}{@{}lX@{}}
\toprule
Documento & Contenido \\
\midrule
\texttt{README.md} & Resumen del reto y del problema con nuestros datos. \\
\texttt{analisis/serie-armonizada.md} & Método de armonización y su validación. \\
\texttt{Documentacion/avance.tex} & Estado del proyecto a la fecha. \\
\texttt{Documentacion/frontend/plan.tex} & Propuesta inicial del frontend. \\
\bottomrule
\end{tabularx}

\section{Contexto: el reto}

Reto \emph{Harmonization of MODIS and VIIRS Hot Spots} del NASA Space Apps Challenge
2026, propuesto por la CONAE (Comisión Nacional de Actividades Espaciales, Argentina).

Los focos de calor que publica NASA FIRMS vienen de dos familias de sensores: MODIS
(desde 2000) y VIIRS (desde 2012). VIIRS detecta muchos más focos que MODIS sobre el
mismo territorio, así que una serie que empalma los dos muestra un aumento de incendios
que no existe. El reto pide resolver esa incompatibilidad y ofrecer el resultado de
forma útil.

\section{Objetivo}

Construir un \textbf{calendario de actividad de quemas} consistente desde el año 2000
hasta hoy, en el que el número de focos signifique lo mismo sin importar el satélite que
los detectó, y ofrecerlo en una \textbf{aplicación web} con la que una persona no
experta pueda responder en menos de un minuto:

\begin{itemize}
  \item ¿Este mes es peor de lo normal en mi zona?
  \item ¿Cuál fue el peor año, y en qué meses se concentra la actividad?
  \item ¿Cuántos focos hubo en esta área y periodo? (y descargarlos)
\end{itemize}

\section{Usuarios}

\begin{tabularx}{\linewidth}{@{}lX@{}}
\toprule
Usuario & Qué necesita \\
\midrule
Gestores de emergencias y bomberos & Saber rápido si el periodo actual es anómalo. \\
Gestores del territorio & Comparar zonas y épocas; planificar temporadas. \\
Científicos & Series consistentes, incertidumbre explícita y descarga de datos. \\
Jurado del reto & Entender el problema y la solución en pocos minutos. \\
\bottomrule
\end{tabularx}

\section{Alcance}

\textbf{Dentro del MVP:}
\begin{itemize}
  \item Armonización MODIS $\to$ VIIRS para 2000--2025 con validación documentada.
  \item Aplicación web con mapa, selección de área y periodo, calendario
        años~$\times$~meses, comparación con lo normal, rankings y descarga en CSV.
  \item Una región de trabajo (India, porque ya tenemos sus datos).
\end{itemize}

\textbf{Fuera del MVP:}
\begin{itemize}
  \item Predicción de incendios.
  \item Datos casi en tiempo real (NRT). Deseable, pero secundario.
  \item Cuentas de usuario.
\end{itemize}

\abierta{\textbf{cobertura mundial.} El README la deja fuera del MVP, pero se ha
planteado mostrar datos de todo el mundo. Implica descargar el archivo global de FIRMS
($\sim$7\,GB comprimidos), recalcular la constante de conversión por región (la actual
solo vale para India) y servir los datos de otra forma (sección~\ref{sec:decisiones}).
Los requerimientos se escriben de modo que la región sea configurable
(RF-01), para que esta decisión no obligue a rehacer la aplicación.}

\section{Convenciones}

Cada requerimiento tiene un identificador y una prioridad:

\begin{tabularx}{\linewidth}{@{}lX@{}}
\toprule
Prefijo & Tipo \\
\midrule
\texttt{RD} & Requerimiento de datos (capítulo~\ref{cap:datos}). \\
\texttt{RF} & Requerimiento funcional de la aplicación (capítulo~\ref{cap:funcionales}). \\
\texttt{RNF} & Requerimiento no funcional o de calidad (capítulo~\ref{cap:nofuncionales}). \\
\bottomrule
\end{tabularx}

\begin{tabularx}{\linewidth}{@{}lX@{}}
\toprule
Prioridad & Significado \\
\midrule
\imp & Sin esto no hay entrega. \\
\rel & Debe estar en la entrega; se recorta solo si no hay tiempo. \\
\des & Mejora si sobra tiempo. \\
\bottomrule
\end{tabularx}

\section{Glosario}

\begin{description}[style=nextline]
  \item[Foco de calor] Píxel en el que un satélite detecta una anomalía térmica.
        Normalmente un fuego activo, pero también puede ser una industria o un volcán.
  \item[FIRMS] \emph{Fire Information for Resource Management System}, el servicio de
        NASA que distribuye los focos de MODIS y VIIRS.
  \item[MODIS] Sensor de los satélites Terra (desde noviembre de 2000) y Aqua (desde
        julio de 2002). Píxel de $\sim$1\,km.
  \item[VIIRS] Sensor de S-NPP (2012), NOAA-20 (2018) y NOAA-21 (2023). Píxel de
        $\sim$375\,m; detecta fuegos más pequeños.
  \item[Armonizar] Llevar los focos de los dos sensores a una escala común. Aquí, la de
        VIIRS S-NPP.
  \item[Constante $k$] Factor por el que se multiplica cada foco MODIS para expresarlo
        en focos VIIRS equivalentes. Se calcula por mes del año.
  \item[Peso] Valor de cada foco en la escala armonizada: $k$ para MODIS, 1 para VIIRS.
        Los focos armonizados de un área son la suma de los pesos.
  \item[Lo normal] Valor de referencia de un mes o año para un área (mediana o promedio
        histórico; decisión abierta).
  \item[Anomalía] Diferencia entre un periodo y lo normal, en porcentaje.
\end{description}

\chapter{El problema}\label{cap:problema}

\section{Datos disponibles}

CSV anuales de FIRMS para India (\texttt{../data de prueba/}, fuera del repo):

\begin{tabularx}{\linewidth}{@{}lllX@{}}
\toprule
Sensor & Archivo & Años & Nota \\
\midrule
MODIS (Terra + Aqua) & \texttt{modis} & 2000--2025 & 2000 solo nov.--dic.; Aqua desde jul.\ 2002. \\
VIIRS S-NPP & \texttt{viirs-snpp} & 2012--2025 & Escala de referencia. \\
VIIRS NOAA-20 & \texttt{viirs-jpss1} & 2018--2025 & Desde abril de 2018. \\
\bottomrule
\end{tabularx}

Se descargan sin clave desde
\url{https://firms.modaps.eosdis.nasa.gov/data/country/<sensor>/<año>/<sensor>_<año>_India.csv}.
El mismo servidor ofrece el archivo de todos los países por año
(\url{data/country/zips/<sensor>_<año>_all_countries.zip}).

\section{La diferencia entre sensores}

\begin{tabularx}{\linewidth}{@{}lXX@{}}
\toprule
& MODIS & VIIRS 375\,m \\
\midrule
Píxel & $\sim$1\,km & $\sim$375\,m \\
Sensibilidad & Fuegos grandes o calientes & También fuegos pequeños \\
\texttt{confidence} & Numérica, 0--100 & Categórica: \texttt{l} / \texttt{n} / \texttt{h} \\
Brillo & \texttt{brightness}, \texttt{bright\_t31} & \texttt{bright\_ti4}, \texttt{bright\_ti5} \\
\texttt{satellite} & \texttt{Terra}, \texttt{Aqua} & \texttt{N} (S-NPP), \texttt{N20} (NOAA-20) \\
\bottomrule
\end{tabularx}

Focos por año en India:

\begin{tabularx}{\linewidth}{@{}XXXX@{}}
\toprule
Año & MODIS & VIIRS S-NPP & S-NPP / MODIS \\
\midrule
2012 & 93\,150 & 589\,159 & 6,3$\times$ \\
2016 & 88\,809 & 681\,943 & 7,7$\times$ \\
2020 & 76\,021 & 489\,632 & 6,4$\times$ \\
2024 & 74\,029 & 552\,312 & 7,5$\times$ \\
\bottomrule
\end{tabularx}

\textbf{Consecuencia.} Si se usa MODIS hasta 2011 y VIIRS después, el calendario muestra
en 2012 un aumento de $\times$5--7 que no existe. Además la proporción no es constante:
cambia con el año, el mes y probablemente la región, así que no basta con un único
factor.

\section{Otros problemas de los datos}

\begin{itemize}
  \item \textbf{Años incompletos:} 2000 solo tiene noviembre y diciembre; Aqua entra en
        julio de 2002. Entre 2000 y 2002 solo hay un satélite MODIS.
  \item \textbf{Tipo de foco:} \texttt{type} = 0 vegetación, 2 fuente estática,
        3 costa afuera. Para quemas solo sirve \texttt{type} = 0.
  \item \textbf{Versión de MODIS:} pasa de 6.2 a 6.03 en 2018; hay que confirmar que no
        introduce un salto.
  \item \textbf{Volumen:} India son $\sim$6,6 millones de focos filtrados en 2000--2025.
        Todo el mundo es del orden de cientos de millones, lo que condiciona cómo se
        sirven los datos a la aplicación.
\end{itemize}

\section{Estado de la solución}

La armonización para India ya está hecha y validada (detalle en
\path{analisis/serie-armonizada.md}):

\begin{itemize}
  \item Salto en 2012 (media 2012--14 / 2009--11): \textbf{5,36} en el empalme crudo,
        \textbf{0,96} en la serie armonizada (1,00 sería no tener salto).
  \item MODIS$\times k$ frente a VIIRS observado, por mes: $r = 0{,}991$ y sesgo 0\,\%
        en calibración (2012--2018); $r = 0{,}987$ y sesgo +2,7\,\% en prueba (2019--2025).
\end{itemize}

\chapter{Requerimientos de datos}\label{cap:datos}

Qué tiene que garantizar la parte de datos: de dónde salen los focos, cómo se filtran,
cómo se armonizan y qué productos entrega a la aplicación.

\section{Fuentes y carga}

\begin{requerimientos}
RD-01 & Los datos se obtienen de los archivos públicos de NASA FIRMS (MODIS colección 6.1
        y VIIRS S-NPP 375\,m), sin consultas en vivo desde la aplicación. & \imp \\
RD-02 & La carga es reproducible con un solo comando (\texttt{mise run datos:cargar})
        a partir de los CSV descargados. & \imp \\
RD-03 & La región de trabajo es un parámetro: cambiar de India a otra región o al mundo
        no requiere reescribir los scripts. & \rel \\
RD-04 & Los datos crudos y la base de datos local no se suben al repositorio; sí los
        scripts que los generan y la URL de descarga. & \imp \\
\end{requerimientos}

\section{Filtrado}

\begin{requerimientos}
RD-05 & Solo se usan focos de vegetación (\texttt{type} = 0). & \imp \\
RD-06 & MODIS: \texttt{confidence} $\geq$ 30. VIIRS: \texttt{confidence} en
        \texttt{n} o \texttt{h}. & \imp \\
RD-07 & Los mismos filtros se aplican en la calibración, en la armonización y en lo que
        muestra la aplicación. & \imp \\
\end{requerimientos}

\section{Armonización}

\begin{requerimientos}
RD-08 & La escala común es la de VIIRS S-NPP: desde febrero de 2012 cada foco S-NPP
        pesa 1. & \imp \\
RD-09 & Entre 2002-08 y 2012-01 cada foco MODIS (Terra + Aqua) pesa $k_\text{modis}$ del
        mes del año; entre 2000-11 y 2002-07, $k_\text{terra}$ del mes del año. & \imp \\
RD-10 & Las constantes se calibran con el solape 2012-02 a 2018-12 y se validan con
        2019--2025, que no se usa para calibrar. & \imp \\
RD-11 & La armonización se hace por foco (columna \texttt{peso}), no por total, para que
        cualquier área se pueda sumar sin recalcular. & \imp \\
RD-12 & Si la región de trabajo es mayor que India, las constantes se calculan por
        región (bioma o franja de latitud) además de por mes del año. & \rel \\
RD-13 & Evaluar si incorporar NOAA-20 (2018--) a la escala mejora la serie, o si basta
        con S-NPP. & \des \\
\end{requerimientos}

\section{Validación e incertidumbre}

\begin{requerimientos}
RD-14 & Se documentan correlación, sesgo, error mediano y RMSE de MODIS$\times k$ frente
        a VIIRS, en calibración y en prueba. & \imp \\
RD-15 & La serie armonizada no muestra salto en 2012: la razón entre la media de
        2012--14 y la de 2009--11 está entre 0,9 y 1,1. & \imp \\
RD-16 & Cada mes armonizado a partir de MODIS lleva una banda de incertidumbre (a partir
        del rango intercuartil de $k$). & \rel \\
RD-17 & Los meses de 2000--2002 (solo Terra) quedan marcados como de menor
        precisión. & \imp \\
\end{requerimientos}

\section{Productos para la aplicación}\label{sec:contrato}

Es el contrato entre la parte de datos y el frontend. Se acuerda antes de construir
la pantalla correspondiente.

\begin{requerimientos}
RD-18 & \textbf{Celdas por mes:} rejilla regular con, por celda y mes, focos
        armonizados ($\sum$\texttt{peso}), focos brutos y fuente
        (\texttt{TERRA}/\texttt{MODIS}/\texttt{VIIRS}). & \imp \\
RD-19 & \textbf{Constantes:} $k$ por mes del año (y región) con su rango
        intercuartil. & \rel \\
RD-20 & \textbf{Focos individuales} con su peso, particionados por año, para la descarga
        en CSV. & \rel \\
RD-21 & \textbf{Metadatos:} región, periodo cubierto, filtros aplicados y versión del
        método. & \rel \\
RD-22 & Todos los productos se generan con tareas de \texttt{mise} desde la base de
        datos local (p.\,ej. \texttt{mise run datos:mapa}). & \imp \\
\end{requerimientos}

Estado: hoy existe un adelanto de RD-18 en JSON con celdas de $0{,}1^\circ$
(\path{analisis/05_exportar_mapa.sql}, 16\,MB para India).

\abierta{tamaño de celda ($0{,}1^\circ$ frente a $0{,}25^\circ$) y formato (JSON,
Parquet o teselas vectoriales). Para el mundo, $0{,}1^\circ$ en JSON no es viable.}

\chapter{Requerimientos funcionales}\label{cap:funcionales}

Qué puede hacer una persona con la aplicación web. Se agrupan por la parte de la
pantalla a la que pertenecen.

\section{Mapa y selección de área}

\begin{requerimientos}
RF-01 & La aplicación muestra un mapa de la región de trabajo con los focos armonizados
        agregados en celdas y coloreados por cantidad. La región es configurable. & \imp \\
RF-02 & El usuario puede elegir un área dibujando un rectángulo en el mapa. & \imp \\
RF-03 & El usuario puede elegir un área escribiendo latitud, longitud y radio. & \rel \\
RF-04 & El usuario puede dibujar un polígono libre. & \des \\
RF-05 & Existe la opción ``toda la región'' como área por defecto. & \imp \\
RF-06 & Al pasar el cursor o tocar una celda se ven sus focos armonizados y
        brutos. & \rel \\
\end{requerimientos}

\section{Selección de periodo}

\begin{requerimientos}
RF-07 & El usuario elige un año (2000--2025) y un mes o el año completo; el mapa se
        actualiza. & \imp \\
RF-08 & El usuario puede elegir un rango de años o de meses. & \rel \\
RF-09 & El usuario puede consultar un día concreto. & \des \\
\end{requerimientos}

\section{Calendario}

\begin{requerimientos}
RF-10 & Para el área elegida, la aplicación muestra un calendario de años (filas,
        2000--2025) por meses (columnas), con cada celda coloreada por focos
        armonizados. Es la vista central. & \imp \\
RF-11 & El calendario se puede ver en valores absolutos o en anomalía respecto a lo
        normal. & \rel \\
RF-12 & Las celdas de 2000--2002 y las de pocos focos se distinguen como de menor
        precisión, sin ocultarlas. & \imp \\
RF-13 & Al pasar el cursor o tocar una celda del calendario se ven: focos armonizados,
        focos brutos, sensor de origen y banda de incertidumbre. & \rel \\
RF-14 & Elegir una celda del calendario lleva el mapa a ese mes. & \des \\
\end{requerimientos}

\section{Comparación con lo normal}

\begin{requerimientos}
RF-15 & Para el periodo elegido se muestra una frase de respuesta, p.\,ej.
        ``Agosto de 2025: 38\,\% por encima de lo normal --- el 3.\textsuperscript{er}
        agosto más activo desde 2000''. & \imp \\
RF-16 & Se muestran lo normal (mediana o promedio histórico), la anomalía en \% y la
        posición en el ranking. & \imp \\
RF-17 & Se listan los 5 meses y los 5 años con más actividad del área. & \rel \\
\end{requerimientos}

\section{Descarga}

\begin{requerimientos}
RF-18 & El usuario descarga en CSV los focos del área y periodo elegidos, con su
        peso. & \imp \\
RF-19 & El usuario descarga el calendario como CSV. & \rel \\
RF-20 & El usuario descarga el calendario como imagen PNG. & \des \\
RF-21 & Los totales del CSV coinciden con lo que muestran el mapa y el
        calendario. & \imp \\
\end{requerimientos}

\section{Ayuda y contexto}

\begin{requerimientos}
RF-22 & La aplicación explica en pocas líneas qué es un foco armonizado y por qué hace
        falta. & \rel \\
RF-23 & Siempre se indica la fuente de los datos (NASA FIRMS) y el método. & \imp \\
\end{requerimientos}

\section{Estado actual}

El prototipo del frontend (rama \texttt{frontend-propuesta}) cubre RF-01 (solo India),
RF-06 y RF-07. El resto está pendiente.

\chapter{Requerimientos no funcionales}\label{cap:nofuncionales}

Cómo de bien tiene que funcionar todo lo anterior.

\section{Rendimiento}

\begin{requerimientos}
RNF-01 & Cualquier consulta de área y periodo responde en menos de 1\,s. & \imp \\
RNF-02 & La primera carga de la aplicación transfiere menos de 5\,MB antes de mostrar
         el mapa. & \rel \\
RNF-03 & Los datos están precalculados; la aplicación no consulta FIRMS en vivo. & \imp \\
\end{requerimientos}

\section{Usabilidad}

\begin{requerimientos}
RNF-04 & Un usuario nuevo responde ``¿este mes es peor de lo normal?'' en menos de un
         minuto, sin ayuda. & \imp \\
RNF-05 & Una sola pantalla principal; nada de menús profundos ni pestañas
         anidadas. & \rel \\
RNF-06 & Diseño sobrio: fondo neutro, un único color de acento; el color se reserva para
         codificar datos. & \rel \\
RNF-07 & Funciona en móvil (pantallas desde 360\,px de ancho): el mapa y el calendario
         se leen sin desplazamiento horizontal. & \rel \\
\end{requerimientos}

\section{Accesibilidad}

\begin{requerimientos}
RNF-08 & Las rampas de color son perceptualmente uniformes y se distinguen con los tres
         tipos comunes de daltonismo (secuencial tipo YlOrRd o viridis; divergente con el
         centro en ``normal'' para anomalías). & \imp \\
RNF-09 & El texto cumple contraste WCAG AA. & \rel \\
RNF-10 & Los controles se pueden usar con teclado y tienen etiquetas para lectores de
         pantalla. & \rel \\
\end{requerimientos}

\section{Idioma}

\begin{requerimientos}
RNF-11 & La interfaz está en español. & \imp \\
RNF-12 & La interfaz está también en inglés para el jurado. & \des \\
\end{requerimientos}

\section{Portabilidad y entorno}

\begin{requerimientos}
RNF-13 & Las herramientas y tareas del proyecto se gestionan con \texttt{mise}
         (\texttt{mise run <tarea>}). & \imp \\
RNF-14 & El entorno de desarrollo arranca con Docker Compose en Docker Desktop
         (Windows), Docker en Linux y Podman. & \imp \\
RNF-15 & La aplicación se puede desplegar como sitio estático (GitHub Pages o
         similar). & \rel \\
RNF-16 & La aplicación se puede abrir desde otros dispositivos de la red (p.\,ej.\
         Tailscale) durante el desarrollo, para probarla en móvil. & \des \\
\end{requerimientos}

\section{Mantenibilidad y reproducibilidad}

\begin{requerimientos}
RNF-17 & Todo resultado (tablas, cifras de validación, archivos de la aplicación) se
         regenera desde los datos crudos con tareas de \texttt{mise}. & \imp \\
RNF-18 & El código pasa el linter (\texttt{mise run front:lint}) y compila sin errores de
         tipos. & \rel \\
RNF-19 & La documentación se escribe en español: \texttt{.md} para el análisis y
         \texttt{.tex} en \texttt{Documentacion/} para planes y requerimientos. & \rel \\
RNF-20 & Los archivos de texto usan fin de línea LF (\texttt{.gitattributes}). & \rel \\
\end{requerimientos}

\section{Licencias y atribución}

\begin{requerimientos}
RNF-21 & Se cita a NASA FIRMS como fuente de los datos y a OpenStreetMap / OpenFreeMap
         como mapa base. & \imp \\
RNF-22 & Solo se usan servicios y librerías gratuitos y sin clave de API. & \rel \\
\end{requerimientos}

\chapter{Arquitectura y estructura del proyecto}\label{cap:arquitectura}

\section{Visión general}

El proyecto tiene dos partes que se comunican por archivos:

\begin{enumerate}
  \item \textbf{Datos} (\texttt{analisis/}): scripts SQL de DuckDB que cargan los CSV de
        FIRMS, calculan las constantes, armonizan cada foco y exportan los productos de
        la sección~\ref{sec:contrato}.
  \item \textbf{Aplicación} (\texttt{frontend/}): página web que lee esos productos como
        archivos estáticos y hace todos los cálculos de área y periodo en el navegador.
\end{enumerate}

Flujo de los datos:

\begin{center}
\small
CSV de FIRMS $\xrightarrow{\text{01\_cargar}}$
\texttt{focos.duckdb} $\xrightarrow{\text{02, 03}}$
constantes $k$ $\xrightarrow{\text{04\_armonizar}}$
focos con peso $\xrightarrow{\text{05\_exportar}}$
\texttt{frontend/public/datos/} $\to$ navegador
\end{center}

\section{Estructura del repositorio}

\begin{verbatim}
Hots-Spots-Nasa-Challenge/
|-- README.md              Reto, problema y criterios de éxito
|-- mise.toml              Herramientas (versiones) y tareas del proyecto
|-- compose.yaml           Entorno Docker: frontend + backend
|-- analisis/              Parte de datos (DuckDB)
|   |-- 01_cargar.sql          CSV -> tabla focos
|   |-- 02_correlacion.sql     Correlación MODIS-VIIRS
|   |-- 03_k_mensual.sql       Constantes por mes
|   |-- 04_armonizar.sql       Peso por foco y serie mensual
|   |-- 05_exportar_mapa.sql   Rejilla 0,1° x mes para el mapa
|   |-- *.md                   Resultados y validación
|   `-- focos.duckdb           Base local (no se versiona)
|-- frontend/              Aplicación web
|   |-- src/
|   |   |-- App.tsx            Pantalla principal y controles
|   |   |-- Mapa.tsx           Mapa (MapLibre GL)
|   |   `-- datos.ts           Carga y agregación de los datos
|   `-- public/datos/      Productos exportados (JSON)
|-- backend/               API FastAPI (esqueleto, solo /salud)
`-- Documentacion/         Este documento y los planes
    |-- main.tex               Documento de requerimientos
    |-- NN-capitulo/           Un capítulo por subcarpeta
    |-- avance.tex             Estado del proyecto
    `-- frontend/plan.tex      Propuesta inicial del frontend
\end{verbatim}

Los CSV crudos viven fuera del repo, en \texttt{../data de prueba/}.

\section{Tecnologías}

\begin{tabularx}{\linewidth}{@{}lX@{}}
\toprule
Pieza & Elección y motivo \\
\midrule
Datos & DuckDB: SQL sobre millones de filas en un portátil, sin servidor. \\
Frontend & React 19 + TypeScript + Vite + Tailwind 4. \\
Mapa & MapLibre GL JS con el mapa base \emph{Positron} de OpenFreeMap (claro, sin
       clave de API). \\
Consultas en el navegador & Hoy, JSON agregado en JavaScript. Previsto: Parquet +
       DuckDB-WASM para áreas y periodos arbitrarios. \\
Gráficos & SVG propio para el calendario (es una cuadrícula). \\
Backend & FastAPI + DuckDB, solo si hace falta (ver decisiones abiertas). \\
Entorno & \texttt{mise} (herramientas y tareas), Docker Compose / Podman,
          \texttt{uv} para Python, \texttt{pnpm} para Node, Tectonic para \LaTeX. \\
\bottomrule
\end{tabularx}

\section{Pantalla principal}

Una sola pantalla, de arriba abajo:

\begin{enumerate}
  \item Barra superior mínima: nombre, región, ayuda.
  \item Selector de periodo (año y mes).
  \item Mapa con la herramienta de selección de área.
  \item Frase de respuesta (anomalía y ranking del periodo).
  \item Calendario años~$\times$~meses con su leyenda.
  \item Rankings y descarga, discretos, al final.
\end{enumerate}

\abierta{si el mapa ocupa la mayor parte de la pantalla con el calendario como panel
lateral, o si el calendario es el protagonista y el mapa un selector pequeño.}

\section{Tareas de \texttt{mise}}

\begin{tabularx}{\linewidth}{@{}lX@{}}
\toprule
Tarea & Qué hace \\
\midrule
\texttt{datos:cargar} & Carga los CSV en \texttt{analisis/focos.duckdb}. \\
\texttt{datos:correlacion} & Correlación y constantes de conversión. \\
\texttt{datos:armonizar} & Serie armonizada 2000--2025. \\
\texttt{datos:mapa} & Exporta la rejilla para el mapa. \\
\texttt{front:dev} / \texttt{front:build} & Servidor de desarrollo / build de producción. \\
\texttt{back:dev} & API en modo desarrollo. \\
\texttt{docker:up} / \texttt{docker:down} & Entorno completo en contenedores. \\
\texttt{docs:pdf} & Compila los \texttt{.tex} de \texttt{Documentacion/}. \\
\bottomrule
\end{tabularx}

\chapter{Plan, aceptación y decisiones abiertas}\label{cap:plan}

\section{Fases}

\begin{tabularx}{\linewidth}{@{}clXl@{}}
\toprule
Fase & Nombre & Entregable & Estado \\
\midrule
1 & Armonización & Serie 2000--2025 validada (India). & Hecho \\
2 & Requerimientos & Este documento, acordado por el equipo. & En curso \\
3 & Contrato de datos & Productos de la sección~\ref{sec:contrato} con datos reales. & En curso \\
4 & Mapa & Mapa con selector de año y mes (RF-01, RF-06, RF-07). & Prototipo \\
5 & Área y calendario & Selección de área y calendario del área (RF-02, RF-05, RF-10). & Pendiente \\
6 & Comparación & Frase de respuesta, anomalías y rankings (RF-15 a RF-17). & Pendiente \\
7 & Descarga & CSV de focos y del calendario (RF-18, RF-19, RF-21). & Pendiente \\
8 & Pulido & Accesibilidad, móvil, textos de ayuda, prueba del minuto. & Pendiente \\
\bottomrule
\end{tabularx}

\section{Criterios de aceptación}

La entrega se considera completa cuando:

\begin{itemize}
  \item Todos los requerimientos \imp{} están cumplidos.
  \item La serie armonizada no muestra salto en 2012 (RD-15) y las métricas de
        validación están documentadas (RD-14).
  \item Un usuario nuevo responde ``¿este mes es peor de lo normal?'' en menos de un
        minuto (RNF-04).
  \item Cualquier consulta de área y periodo responde en menos de 1\,s (RNF-01).
  \item El CSV descargado coincide con lo que muestran el mapa y el calendario (RF-21).
  \item La rampa de color pasa la prueba de daltonismo (RNF-08).
\end{itemize}

\section{Decisiones abiertas}\label{sec:decisiones}

\begin{enumerate}
  \item \textbf{Región: India o el mundo.} El README fija India para el MVP; se ha
        propuesto mostrar datos mundiales. El mundo implica:
        \begin{itemize}
          \item descargar $\sim$7\,GB comprimidos (MODIS 2000--2025 y VIIRS S-NPP
                2012--2025, archivo \emph{all countries} de FIRMS);
          \item recalcular $k$ por región, porque la de India no vale para otros biomas
                ni para el hemisferio sur (RD-12);
          \item cambiar el formato de los datos de la aplicación (Parquet o teselas
                vectoriales en vez de JSON).
        \end{itemize}
        Opción intermedia: India para el MVP y el mundo como siguiente paso, con la
        región ya configurable (RD-03, RF-01).
  \item \textbf{Región de la demo:} India (datos ya disponibles) o Argentina (región de
        la CONAE, que propone el reto).
  \item \textbf{Lo normal:} mediana o promedio; sobre todo 2000--2025 o sobre una
        ventana móvil (p.\,ej.\ los 10 años anteriores).
  \item \textbf{Backend o estático:} el plan inicial era sin servidor, pero existe un
        esqueleto de API. Depende de si la descarga de focos individuales cabe en el
        navegador.
  \item \textbf{Celda y formato} de los datos del mapa ($0{,}1^\circ$ o $0{,}25^\circ$;
        JSON, Parquet o teselas).
  \item \textbf{Mostrar números absolutos o un índice relativo}, para no dar una falsa
        sensación de precisión.
  \item \textbf{Idioma:} solo español, o también inglés para el jurado.
  \item \textbf{¿Subir los datos exportados al repositorio?} Sin ellos, quien no tenga
        los CSV ve la aplicación vacía.
\end{enumerate}

\section{Historial de versiones}

\begin{tabularx}{\linewidth}{@{}llX@{}}
\toprule
Versión & Fecha & Cambios \\
\midrule
0.1 & 2026-11-14 & Primer borrador: requerimientos de datos, funcionales y no
                   funcionales; estructura del proyecto; decisiones abiertas. \\
\bottomrule
\end{tabularx}


\end{document}
